
\section{The regulated BV theories in \texorpdfstring{$D$}{D} dimensions}
\label{sec:schemes}

A dimensional scheme is a BV theory in $D$ dimensions: a field space, an odd pairing, a bracket and a regulated equivariant action whose failure to satisfy \eqref{eq:MC4} is a local cochain that vanishes in the four-dimensional limit. The Clifford algebra in $D$ dimensions exists only abstractly, and the two schemes differ in how much of the four- and ten-dimensional spinor structure survives.

\subsection{Formal spaces and the four-dimensional limit}

\begin{axiom}[Formal spaces, {[S05, §2, App.~A], [G17, (2.3)], [BM77]}]
\label{ax:spaces}
There are infinite-dimensional real inner-product spaces with formal dimensions
\[
4S\subset QDS\subset Q4S,\qquad QDS=4S\oplus Q(D{-}4)S,\qquad Q4S=QDS\oplus Q(4{-}D)S,
\]
with $4S$ strictly four-dimensional, such that the orthogonal projections onto $Q4S$, $QDS$, $4S$, $Q(D{-}4)S$, $Q(4{-}D)S$ have traces $4$, $D$, $4$, $D-4$, $4-D$, and a translation-invariant integral $\int\frac{d^Dp}{(2\pi)^D}$ on $QDS$ with $\int d^Dp\,f(sp)=s^{-D}\int d^Dp\,f(p)$ and $\int\frac{d^Dp}{(2\pi)^D}e^{-t|p|^2}=(4\pi t)^{-D/2}$. Integrals of polynomials vanish. Only the inner product and these traces may be used: in $Q4S$ five vectors need not be linearly dependent [S05, §2].
\end{axiom}

\begin{definition}[Evanescent quantities and $\mathrm{LIM}$]
\label{def:lim}
A $D$-dependent expression is in normal form when all Clifford products are written in the antisymmetrised basis $\rho\circ q$, all contractions are carried out with the traces of Axiom \ref{ax:spaces}, and no identity of a finite-dimensional space other than $4S$ and $W^\perp_{\mathbb R}$ is used. For an expression whose Laurent expansion in $4-D$ has no pole, $\mathrm{LIM}$ takes the coefficient of $(4-D)^0$ of its normal form and then applies the strictly four-dimensional (and, in Scheme D, ten-dimensional) identities: components along $Q(D{-}4)S$ and $Q(4{-}D)S$ are set to $0$, Fierz identities and $[\mathrm{ESW},\ \mathrm{Assumption}\ 3.1]$ become available [BIMS20, §2]. A local functional $E$ is evanescent when $\mathrm{LIM}\,E=0$.
\end{definition}

\subsection{Scheme D: dimensional reduction from ten dimensions}

\begin{definition}[Scheme D]
\label{def:schemeD}
\mbox{}\begin{itemize}
\item Spaces: $V_{10}\mapsto V_{10}^{\rm D}:=Q4S\oplus W^\perp_{\mathbb R}$ (complexified), of formal dimension $10$; positions $x\in QDS$; every field is constant along $Q(4{-}D)S\oplus W^\perp_{\mathbb R}$; $d$ acts along $QDS$; $\mathrm{dvol}_D$ is the measure of $QDS$.
\item Fields: $A\in\Omega^0(QDS;(V_{10}^{\rm D})^\vee\otimes\mathfrak g)$, with $A|_{QDS}$ the connection, $A|_{Q(4-D)S}$ the $4-D$ $\epsilon$-scalars and $A|_{W^\perp_{\mathbb R}}$ the six scalars: a $D$-dimensional gauge field and $10-D$ scalars; $\lambda\in\Omega^0(QDS;\Pi\Sigma_D\otimes\mathfrak g)$; $c$, antifields and degrees as in [HK, def:c1-fields].
\item Clifford data: $\mathrm{Cl}_D$ is the algebra generated by $\rho(v)$, $v\in V_{10}^{\rm D}$, with the single relation $\rho(v_1)\rho(v_2)+\rho(v_2)\rho(v_1)=2(v_1,v_2)$; $\Sigma_D\oplus\Sigma_D^*$ is a $\mathbb Z/2$-graded module on which every $\rho(v)$ is odd (a totally anticommuting chirality, [S05, (2.6)]), with an even pairing $(-,-)\colon\Sigma_D\otimes\Sigma_D^*\to\mathbb C$, $(\rho(v)\zeta_1,\zeta_2)=(\zeta_1,\rho(v)\zeta_2)$, and $\Gamma$ defined by $(\rho(v)\zeta_1,\zeta_2)=(v,\Gamma(\zeta_1,\zeta_2))$. Traces of even products are fixed by $\operatorname{tr}_{\Sigma_D}1=16$, cyclicity and the Clifford relation; the chirality-odd part of every trace is $0$ [S05, footnote 1].
\item Action: $S_D$ is \eqref{eq:S4} with $\mathrm{dvol}\mapsto\mathrm{dvol}_D$, $F_A=dA+\frac12[A\wedge A]$, ${\not d}_A=\rho\circ(d+[A,-])$ with all components of $A$, and the Hodge pairing of $\wedge^2(V_{10}^{\rm D})^\vee$.
\item Ghost spinors: $\Sigma$, $\Sigma^*$ of $\mathfrak{psl}(4|4)$ embed in $\Sigma_D$, $\Sigma_D^*$ as $\mathrm{Cl}(4S\oplus W_{\mathbb R}^\perp)$-modules with $\Gamma(\zeta_1,\zeta_2)\in 4S\oplus W^\perp_{\mathbb R}$ for $\zeta_1,\zeta_2$ in the image.
\end{itemize}
\end{definition}

\begin{remark}[Two-component form of Scheme D]
\label{rem:D2c}
Two-component formulas keep the form of [HK, def:c8-twocomp] with $m\in Q4S$ formal: $\sigma^\mu\bar\sigma^\nu+\sigma^\nu\bar\sigma^\mu=-2g^{\mu\nu}$, $g^\mu{}_\mu=4$, and no further identity. In particular the completeness relation $\sigma^m_{\alpha\dot\alpha}\bar\sigma^{m\dot\beta\beta}=-2\delta_\alpha{}^\beta\delta^{\dot\beta}{}_{\dot\alpha}$, the Schouten identity and the duality $\sigma^{[m}\bar\sigma^n\sigma^{p]}\propto\epsilon^{mnpq}\sigma^q$ are unavailable: they are four-dimensional Fierz identities. The R-symmetry structure is explicit: $\Sigma^I_{ij}$, $\sum_I\Sigma^I_{ij}\Sigma^I_{kl}=2\varepsilon_{ijkl}$, $(\rho(w)\xi)^\alpha{}_j=\varepsilon^{\alpha\beta}w_{ji}\xi_\beta{}^i$ as in [HK, def:c8-twocomp]. The kinetic terms read
\[
S_D\supset\int d^Dx\Big\{-\tfrac18(F_{\mu\nu},F_{\mu\nu})-\tfrac1{16}(D_mA^{ij},D_mA_{ij})+i\big(\bar\lambda_{\dot\alpha i},\bar\sigma^{m\dot\alpha\alpha}D_m\lambda_\alpha{}^i\big)+(D_mc,A^*_m)\Big\},
\]
$\mu,\nu\in Q4S$, $m\in QDS$,
with $F_{m\tilde n}=D_mA_{\tilde n}$ and $F_{\tilde m\tilde n}=[A_{\tilde m},A_{\tilde n}]$ for $\tilde m,\tilde n\in Q(4{-}D)S$; for $\mathfrak{su}(N)$, $S_D/\hbar\supset\int d^Dx\{-\frac14F^a_{\mu\nu}F^a_{\mu\nu}-\frac12(D_mA_I)^a(D_mA_I)^a+i\bar\lambda^a\bar\sigma^mD_m\lambda^a\}$ [HK, prop:c8-twocomp].
\end{remark}

\subsection{Scheme B: Breitenlohner--Maison/'t Hooft--Veltman with an explicit evanescent split}

\begin{proposition}[A constraint on BMHV kinetic operators]
\label{prop:nogo}
Let $V_{10}\mapsto QDS\oplus W^\perp_{\mathbb R}$ with the four-dimensional chirality $\rho(e_1e_2e_3e_4)$ on $4S$ as in [HK, prop:c1-chirality], so that $\rho(\hat v)$ commutes with the ten-dimensional chirality for $\hat v\in Q(D{-}4)S$ and maps $\Sigma\to\Sigma$. A quadratic form $\frac12\int(\lambda,K\lambda)$ on $\Pi\Sigma\otimes\mathfrak g$ alone whose evanescent part is first order in $d|_{Q(D-4)S}$ and preserves the four-dimensional chirality requires an invariant linear map $J\colon W\to W^*$ of the R-symmetry group it preserves. No such map is $SU(4)$-invariant.
\end{proposition}

\begin{proof}
The pairing of [HK, def:c1-spaces] pairs $\Sigma$ with $\Sigma^*$. An evanescent term $(\lambda,\rho(\hat v)\lambda)$ pairs $\Sigma$ with $\rho(\hat v)\Sigma\subset\Sigma$, which needs a map $\Sigma\to\Sigma^*$. Preserving the four-dimensional chirality, it maps $S_+\otimes W\to S_+\otimes W^*$ and $S_-\otimes W^*\to S_-\otimes W$, so it is $1_{S}\otimes J$ with $J\colon W\to W^*$ intertwining the R-symmetry. $W$ is the fundamental representation of $SU(4)$ and $W^*$ the antifundamental; they are inequivalent, so $J=0$ for $SU(4)$. Adding a second spinor field in $\Pi\Sigma^*\otimes\mathfrak g$ with a gauge-covariant four-dimensional kinetic term changes the four-dimensional field content.
\end{proof}

\begin{definition}[Scheme B]
\label{def:schemeB}
Fix the real structure $J=\delta\colon W\to W^*$, $J(z)_i=\delta_{ij}z^j$, preserved by $SO(4)_R\subset SU(4)_R$.
\begin{itemize}
\item Spaces: $V_{10}\mapsto V_{10}^{\rm B}:=QDS\oplus W^\perp_{\mathbb R}$, $QDS=4S\oplus Q(D{-}4)S$; positions in $QDS$; fields constant along $W^\perp_{\mathbb R}$.
\item Fields: $A\in\Omega^0(QDS;(V_{10}^{\rm B})^\vee\otimes\mathfrak g)$, i.e. a $D$-dimensional gauge field $A|_{QDS}=A|_{4S}+A|_{Q(D-4)S}$ and six scalars $A|_{W^\perp_{\mathbb R}}$ (no $\epsilon$-scalars); $\lambda\in\Omega^0(QDS;\Pi(\Sigma\otimes M)\otimes\mathfrak g)$, where $\Sigma$ is the four-dimensional semi-spin space of [HK] and $M$ is a $\mathbb Z/2$-graded module of $\mathrm{Cl}(Q(D{-}4)S)$ with generators $\tau(\hat v)$, $\tau(\hat v_1)\tau(\hat v_2)+\tau(\hat v_2)\tau(\hat v_1)=-2(\hat v_1,\hat v_2)$, an odd pairing $(m,\tau(\hat v)m')=-(m',\tau(\hat v)m)$ and $\operatorname{tr}_M1=1$.
\item Clifford data: $\rho(v)$ for $v\in 4S\oplus W^\perp_{\mathbb R}$ is the four-dimensional one of [HK] tensored with $1_M$; for $\hat v\in Q(D{-}4)S$, $\rho_J(\hat v):=i\,\gamma_5\,(J\otimes\tau(\hat v))$ on $\Sigma\oplus\Sigma^*$, with $J$ on the factor $W$ and $J^{-1}$ on $W^*$, and $\gamma_5$ the four-dimensional chirality below ($+1$ on $S_+$, $-1$ on $S_-$). Then $\rho(v)$, $\rho_J(\hat v)$ satisfy the Clifford relations of $QDS$ among themselves; with $W^\perp_\pm$ the eigenspaces $\pm1$ of $w_{ij}\mapsto\delta_{ik}\delta_{jl}w^{kl}$ (two $SO(4)_R$-triples), $\rho_J(\hat v)$ anticommutes with $\rho(w)$ for $w$ in one of them and commutes with $\rho(w)$ for $w$ in the other. $\gamma_5:=\rho(e_1e_2e_3e_4)$ is the four-dimensional chirality; it anticommutes with $\rho(4S)$ and commutes with $\rho_J(Q(D{-}4)S)$ [BM77], [BIMS20, §2].
\item Action: $S^{\rm B}_D=S_{4S}+S_{\rm evan}$, where $S_{4S}$ is \eqref{eq:S4} with $\mathrm{dvol}_D$, all contractions over $4S\oplus W^\perp_{\mathbb R}$, the gauge kinetic term over all of $QDS\oplus W^\perp_{\mathbb R}$ on the components without $Q(D{-}4)S$-index, and
\[
S_{\rm evan}=\int\mathrm{dvol}_D\Big(-\tfrac14(F_A,F_A)\big|_{\rm hat}+\tfrac12\big(\lambda,\rho_J(d_A|_{Q(D-4)S})\lambda\big)\Big),
\]
$(F_A,F_A)|_{\rm hat}$ being the terms of $(F_A,F_A)$ with at least one index along $Q(D{-}4)S$. The fermion kinetic operator is $\rho(d|_{4S})+\rho_J(d|_{Q(D-4)S})$, with inverse $\frac{\rho(p|_{4S})+\rho_J(p|_{Q(D-4)S})}{|p|^2}$ up to the factor $i$ and $(-,-)^{-1}$.
\end{itemize}
\end{definition}

\begin{remark}[Two-component and Majorana form of Scheme B]
\label{rem:B2c}
$J$ identifies $\bar\lambda^{\dot\alpha A}:=\varepsilon^{\dot\alpha\dot\beta}\delta^{AB}\bar\lambda_{\dot\beta B}$, and $\Psi^A=(\lambda_\alpha{}^A,\bar\lambda^{\dot\alpha A})$, $A=1,\dots,4$, are four Majorana spinors. With the 't Hooft--Veltman matrices $\gamma^\mu=(\gamma^m,\hat\gamma^{\hat m})$, $\{\gamma^\mu,\gamma^\nu\}=2g^{\mu\nu}$ on $QDS$, $\gamma_5$ built from $4S$, the kinetic term is $\frac i2\bar\Psi^A\gamma^\mu D_\mu\Psi^A$ in Srednicki's conventions with the Wick map of [HK, def:c8-twocomp]; its evanescent part $\frac i2\bar\Psi^A\hat\gamma^{\hat m}D_{\hat m}\Psi^A$ reads in two-component form
\[
\tfrac12\int d^Dx\Big\{\delta_{ij}\varepsilon^{\alpha\beta}\big(\lambda_\alpha{}^i,\tau^{\hat m}D_{\hat m}\lambda_\beta{}^j\big)-\delta^{ij}\varepsilon^{\dot\alpha\dot\beta}\big(\bar\lambda_{\dot\alpha i},\tau^{\hat m}D_{\hat m}\bar\lambda_{\dot\beta j}\big)\Big\},
\]
with $\tau^{\hat m}$ acting on the formal factor $M$; the four-dimensional part is $i(\bar\lambda_{\dot\alpha i},\bar\sigma^{m\dot\alpha\alpha}D_m\lambda_\alpha{}^i)$, $m\in 4S$. The six scalars split under $SO(4)_R\cong SU(2)\times SU(2)$ into the triples $W^\perp_\pm$; the Yukawa couplings are those of [HK] on $4S\oplus W^\perp_{\mathbb R}$, equivalently $\bar\Psi^A(\alpha^I_{AB}+i\gamma_5\beta^I_{AB})[A_I,\Psi^B]$ with the 't Hooft symbols of [HK, def:c8-twocomp]. All $\sigma$-matrices in Scheme B are explicit $4S$ matrices; four-dimensional Fierz identities hold on $4S$-structures, and the formal objects are $\tau^{\hat m}$, $\hat p$, $Q(D{-}4)S$-components of $A$.
\end{remark}

\subsection{The non-minimal sector and gauge fixing}

\begin{definition}[Non-minimal pair]
\label{def:nm}
Add the fields $\bar c\in\Omega^0(QDS;\mathfrak g[-1])$ and $b\in\Omega^0(QDS;\mathfrak g)$ and their antifields $\bar c^*$, $b^*$; the coordinate functionals have the degrees
\[
\bar c:\ -1,\qquad b:\ 0,\qquad \bar c^*:\ 0,\qquad b^*:\ -1,
\]
total parities $\bar c$ odd, $b$ even, $\bar c^*$ even, $b^*$ odd; the even members of the two new pairs are $\bar c^*$ and $b$. The pairing, $K_0$ and the bracket extend by the rules of [HK, def:c1-pairing, def:c1-bracket] with the fibre pairing $(-,-)$ of $\mathfrak g$. The non-minimal action is
\[
S_{\rm nm}:=S_D+\int\mathrm{dvol}_D\,(b,\bar c^*),
\]
and $Q$ extends to the unique operator with $\frac12\langle e,Qe\rangle=\frac12\langle e,Qe\rangle_{\rm min}+\int\mathrm{dvol}_D(b,\bar c^*)$ preserving the pairing; it maps the slot of $b$ to that of $\bar c$ and the slot of $\bar c^*$ to that of $b^*$, and $Q^2=0$.
\end{definition}

\begin{definition}[Gauge fermion and gauge-fixed action]
\label{def:gf}
For $\xi_{\rm GF}\in\mathbb R$,
\[
\Psi:=\int\mathrm{dvol}_D\big(\bar c,\ -d^*A|_{QDS}+\xi_{\rm GF}\,b\big),\qquad -d^*A|_{QDS}=\partial\cdot A|_{QDS},
\]
of degree $-1$. The gauge-fixed equivariant action is $S_{\Psi,\mathfrak{psl}}:=e^{\{\Psi,-\}}S_{{\rm nm},\mathfrak{psl}}$ (Definition \ref{def:equiv}); since $\Psi$ contains no antifield, $\{\Psi,-\}$ lowers the antifield number by one and the series terminates; on functionals it is the substitution $\phi^*\mapsto\phi^*+\delta\Psi/\delta\phi$ of [HK]'s pairing conventions. Eliminating $b$ leaves $-\frac1{4\xi_{\rm GF}}\int\mathrm{dvol}_D(d^*A|_{QDS},d^*A|_{QDS})$, which with the gauge kinetic term $-\frac14\int(dA,dA)$ gives $-\frac14\int\big(A,(d^*d+\xi_{\rm GF}^{-1}dd^*)A\big)$: $\xi_{\rm GF}=1$ is the Feynman gauge, $\langle A_mA_n\rangle=2(-,-)^{-1}\delta_{mn}/|p|^2$, and in $\mathfrak{su}(N)$ components the gauge term is $-\frac{\hbar}{2\xi_{\rm GF}}(\partial\cdot A^a)^2$.
\end{definition}

\begin{lemma}[Gauge fixing is a canonical transformation]
\label{lem:canon}
$\{\Psi,-\}$ is an even derivation of the bracket of degree $0$ commuting with $d_{CE}$. Hence $e^{\{\Psi,-\}}$ preserves the bracket and
\[
d_{CE}S_{\Psi,\mathfrak{psl}}+\tfrac12\{S_{\Psi,\mathfrak{psl}},S_{\Psi,\mathfrak{psl}}\}=e^{\{\Psi,-\}}\Big(d_{CE}S_{{\rm nm},\mathfrak{psl}}+\tfrac12\{S_{{\rm nm},\mathfrak{psl}},S_{{\rm nm},\mathfrak{psl}}\}\Big).
\]
The gauge-fixed operator $Q^{GF}$ of the free part maps antifield slots to field slots, $(Q^{GF})^2=0$, and $[Q,Q^{GF}]=H$ with $H=|p|^2$ on every slot at $\xi_{\rm GF}=1$; the propagator is $P_D=(Q^{GF}\otimes1)K_0/|p|^2$ and $(Q\otimes1+1\otimes Q)P_D=K_0$ as an identity of $D$-dimensional kernels.
\end{lemma}

\begin{proof}
$\Psi$ is odd of degree $-1$; by the Jacobi identity of the odd bracket, $\{\Psi,-\}$ is a derivation of degree $0$ of $\{,\}$ and of the product, even as an operator; it is $C^\bullet(\mathfrak{psl}(4|4))$-linear because $\Psi$ has arity $0$, so it commutes with $d_{CE}$. The exponential of an even derivation of a bracket is an automorphism; applied to the Maurer--Cartan defect it gives the displayed identity. For the free part, $e^{\{\Psi,-\}}$ substitutes $A^*|_{QDS}\mapsto A^*|_{QDS}+\delta\Psi/\delta A$, $\bar c^*\mapsto\bar c^*+\delta\Psi/\delta\bar c$, $b^*\mapsto b^*+\delta\Psi/\delta b$, so the field part of the free gauge-fixed action is the quadratic form of Definition \ref{def:gf} and the ghost term $\int(dc,\delta\Psi/\delta A)$, a multiple of $\int(dc,d\bar c)$; its inverse, written as $Q^{GF}/H$ on the antifield slots, is the propagator, and $Q^{GF}$ reads antifields only, hence squares to zero. The kernel identity is $p$-by-$p$: $Q(p)Q^{GF}(p)+Q^{GF}(p)Q(p)=|p|^2$ on each slot, divided by $|p|^2$.
\end{proof}

\begin{remark}
The heat-kernel track uses the minimal sector and a Feynman-gauge $Q^{GF}$ with $(Q^{GF})^2\neq0$ [HK, ch.~3]; the composite $A|_V$--$c$ lines of [HK, ch.~9--12] are products $P\star X P\star$ with $(P\star)^2\propto(Q^{GF})^2$. With the non-minimal pair $(P_D\star)^2=0$; the corresponding contributions reappear along $\xi,\eta,K_{v_1}$ as the vertices of Proposition \ref{prop:gfvert}.
\end{remark}

\subsection{The regulated equivariant action}

\begin{definition}[Transcription rules and the regulated equivariant action]
\label{def:equiv}
$S_{{\rm nm},\mathfrak{psl}}:=S_{\rm nm}+S^{(1)}_D+S^{(2)}_D\in C^\bullet(\mathfrak{psl}(4|4),\mathscr O_l(\mathscr E_D))$, where the ghost algebra is the four-dimensional $\mathfrak{psl}(4|4)$ itself and $S^{(1)}_D$, $S^{(2)}_D$ are [HK, eq:c2-S1even, eq:c2-S1Q, eq:c2-S1S, eq:c2-S2] transcribed by:
\begin{itemize}
\item[T1] $\mathrm{dvol}\mapsto\mathrm{dvol}_D$, $d$ along $QDS$, all pairings as in the scheme;
\item[T2] Killing data in $QDS$: $P_{v_1}$, $v_1\in 4S$; $M_\omega$, $\omega\in\mathfrak{so}(4S)$; $\Delta=2E$, $E(x)=x\in QDS$; $K_{v_1}$, $v(x)=(x,x)v_1-2(v_1,x)x$, $x\in QDS$; Killing spinors $x\mapsto\xi+\rho(x)\eta$, $x\in QDS$; $R_r$, $r\in\mathfrak{so}(W^\perp_{\mathbb R})$;
\item[T3] Weyl factors $\frac14\operatorname{div}v\mapsto\sigma_v$, $\frac38\operatorname{div}v\mapsto\frac32\sigma_v$: four-dimensional weights $A\colon1$ (all components), $\lambda\colon\frac32$, $c\colon0$, $\bar c\colon2$, $b\colon2$; the antifields then carry $D$ minus these weights. With the internal space $E:=Q(4{-}D)S\oplus W^\perp_{\mathbb R}$ in Scheme D and $E:=W^\perp_{\mathbb R}$ in Scheme B, every occurrence of $A|_{W^\perp_{\mathbb R}}$ in [HK, eq:c2-S1even, eq:c2-S1S] becomes $A|_E$: $v.A=L_v(A|_{QDS})+v(A|_E)+\sigma_vA|_E$, and the last term of $S^{(1)}(\eta)$ is $-(\rho(A|_E)\eta,\lambda^*)$;
\item[T4] in Scheme B every structure built from a ghost ($\Gamma(\zeta,\lambda)$, $\rho(F_A)\zeta$, $\rho(A|_{W^\perp})\eta$) uses components along $4S\oplus W^\perp_{\mathbb R}$ only;
\item[T5] the non-minimal pair: for even $x$ with vector field $v$, $S^{(1)}_D(x)\supset-\int(x.\bar c,\bar c^*)+\int(x.b,b^*)$, $x.\bar c=v(\bar c)+2\sigma_v\bar c$, $x.b=v(b)+2\sigma_vb$; odd $x$ act trivially on $(\bar c,b)$; for odd $x_1,x_2$,
\[
S^{(2)}_D(x_1,x_2)\supset\kappa\int\mathrm{dvol}_D\,\big([x_1,x_2].\bar c,\ b^*\big),
\]
with $\kappa\in\{\pm1\}$ fixed by the arity-two component of the Maurer--Cartan equation restricted to $(\bar c,b)$ (Proposition \ref{prop:gfvert} (iii)).
\end{itemize}
\end{definition}

\begin{proposition}[The gauge-fixing vertices]
\label{prop:gfvert}
Write $x.F:=\{S^{(1)}_D(x),F\}$ for a functional $F$ (for odd $x$ the parity sign of [HK, thm:c8-L (iv)] is $+1$), so that $x.A$ is the coefficient of $A^*$ in $S^{(1)}_D(x)$: $\xi.A=-\Gamma(\xi,\lambda)$, $\eta.A=-\Gamma(\rho(x)\eta,\lambda)$, $v.A=L_vA+\sigma_vA|_{W^\perp}$ in the conventions of [HK, def:c2-S1].
\begin{itemize}
\item[(i)] $S^{(1)}_{\Psi}(x)=S^{(1)}_D(x)+\{\Psi,S^{(1)}_D(x)\}$ with $\{\Psi,S^{(1)}_D(x)\}=-x.\Psi$ for even $x$ and $+x.\Psi$ for odd $x$; $S^{(2)}_\Psi=S^{(2)}_D+\{\Psi,S^{(2)}_D\}$.
\item[(ii)] Odd generators ($\lambda$--$\bar c$ vertices): $\xi.\Psi=-\int\mathrm{dvol}_D\big(\bar c,\partial\cdot(\xi.A)|_{QDS}\big)=-\int\mathrm{dvol}_D\big(d\bar c,\Gamma(\xi,\lambda)|_{QDS}\big)$, and the same with $\rho(x)\eta$ in place of $\xi$.
\item[(iii)] Even generators: with $\partial\cdot(L_vA)=(v+2\sigma_v)\partial\cdot A+(2-D)(d\sigma_v,A)$ for conformal Killing $v$ on $QDS$,
\[
v.\Psi=(4-D)\int\mathrm{dvol}_D\,\sigma_v\big(\bar c,\partial\cdot A+\xi_{\rm GF}b\big)+(2-D)\int\mathrm{dvol}_D\,\big(\bar c,\iota_{\mathrm{grad}\,\sigma_v}A\big).
\]
It vanishes for $P_{v_1}$, $M_\omega$, $R_r$; for $\Delta$ it is $2(4-D)\int(\bar c,\partial\cdot A+\xi_{\rm GF}b)$; for $K_{v_1}$ the last term is the $A|_{QDS}$--$\bar c$ vertex $2(D-2)\int(\bar c,\iota_{v_1}A)$, equal to $4\int(\bar c,\iota_{v_1}A)$ at $D=4$.
\item[(iv)] Arity two ($\bar c\bar c$ vertex): $\{\Psi,\kappa\int([x_1,x_2].\bar c,b^*)\}=\pm\kappa\,\xi_{\rm GF}\int\mathrm{dvol}_D\big([x_1,x_2].\bar c,\bar c\big)$, nonzero for $\xi_{\rm GF}\neq0$ and absent in the Landau gauge.
\end{itemize}
\end{proposition}

\begin{proof}
(i) Graded symmetry $\{f,g\}=(-1)^{|f||g|}\{g,f\}$: $\Psi$ is odd, $S^{(1)}(x)$ has parity $1+|x|$ [HK, def:c2-cochains]; $S^{(1)}$ is linear and $S^{(2)}$ quadratic in antifields not paired with fields of $\Psi$ except $A^*|_{QDS}$, $\bar c^*$, $b^*$, and $\{\Psi,\{\Psi,S^{(1)}\}\}=0$ because $\{\Psi,S^{(1)}\}$ contains no antifield of $A$, $\bar c$, $b$. (ii) $\xi.\bar c=\xi.b=0$; the odd derivation $\xi.$ passes the odd $\bar c$ with the sign $-1$; $\int(\bar c,\partial\cdot X)=-\int(d\bar c,X)$; the second form uses $\xi.A=-\Gamma(\xi,\lambda)$. (iii) For a conformal Killing vector, $\partial_\mu v_\nu+\partial_\nu v_\mu=2\sigma_v\delta_{\mu\nu}$ and $\partial^2v_\nu=(2-D)\partial_\nu\sigma_v$, which give the stated divergence identity; then
\[
\begin{aligned}
v.\Psi={}&\int\big((v+2\sigma_v)\bar c,\partial\cdot A\big)+\int\big(\bar c,(v+2\sigma_v)\partial\cdot A\big)+(2-D)\int(\bar c,\iota_{\mathrm{grad}\,\sigma_v}A)\\
&+\xi_{\rm GF}\Big[\int((v+2\sigma_v)\bar c,b)+\int(\bar c,(v+2\sigma_v)b)\Big],
\end{aligned}
\]
and $\int\big((v+2\sigma_v)f,g\big)+\int\big(f,(v+2\sigma_v)g\big)=\int(4\sigma_v-\operatorname{div}v)(f,g)=(4-D)\int\sigma_v(f,g)$. For $K_{v_1}$, $\mathrm{grad}\,\sigma_v=-2v_1$. (iv) $\delta\Psi/\delta b=\xi_{\rm GF}\bar c$.
\end{proof}

\begin{remark}[Two-component form of the gauge-fixing vertices]
For $\xi=(\xi_\alpha{}^i,\bar\xi_{\dot\alpha i})$, $\Gamma(\xi,\lambda)|_V^m=i\big(\bar\lambda_{\dot\alpha i}\bar\sigma^{m\dot\alpha\alpha}\xi_\alpha{}^i+\bar\xi_{\dot\alpha i}\bar\sigma^{m\dot\alpha\alpha}\lambda_\alpha{}^i\big)$ [HK, def:c8-twocomp], so
\[
\xi.\Psi=-i\int d^Dx\Big(\partial_m\bar c,\ \bar\lambda_{\dot\alpha i}\bar\sigma^{m\dot\alpha\alpha}\xi_\alpha{}^i+\bar\xi_{\dot\alpha i}\bar\sigma^{m\dot\alpha\alpha}\lambda_\alpha{}^i\Big),\qquad m\in QDS\ (\text{Scheme D}),\ m\in4S\ (\text{Scheme B}),
\]
and along $\eta$ the same with $\epsilon_\alpha{}^i(x)=-ix_m\sigma^m_{\alpha\dot\alpha}\bar\eta^{\dot\alpha i}$, $\bar\epsilon_{\dot\alpha i}(x)=ix_m\eta^\alpha{}_i\sigma^m_{\alpha\dot\alpha}$ [HK, eq:c8-killing]; this is the structure $\partial^\mu\bar c(\epsilon\sigma_\mu\bar\lambda+\lambda\sigma_\mu\bar\epsilon)$ of [MPW96, (2.8)]. The $K_{v_1}$ vertex is $2(D-2)\int d^Dx\,(\bar c,v_{1,m}A_m)$. The arity-two vertex along $(\xi,\bar\xi')$, $\xi\in S_+\otimes W$, $\bar\xi'\in S_-\otimes W^*$, is $\pm\kappa\,\xi_{\rm GF}\,i\,\bar\xi'_{\dot\alpha i}\bar\sigma^{m\dot\alpha\alpha}\xi_\alpha{}^i\int d^Dx\,(\partial_m\bar c,\bar c)$, the counterpart of $s\,b\supset-2i\epsilon\sigma^\mu\bar\epsilon\,\partial_\mu\bar c$ in [MPW96, (2.3)].
\end{remark}

\subsection{The classical breaking}

\begin{definition}[Classical breaking]
\label{def:breaking}
$\widehat\Delta:=d_{CE}S_{\Psi,\mathfrak{psl}}+\frac12\{S_{\Psi,\mathfrak{psl}},S_{\Psi,\mathfrak{psl}}\}$, a local cochain of degree $1$, with components $\widehat\Delta^{(k)}(x_1,\dots,x_k)$ normalised as in [HK, lem:c2-components]; $b_D:=d_{CE}+\{S_{\Psi,\mathfrak{psl}},-\}$.
\end{definition}

\begin{lemma}[Bianchi identity of the breaking]
\label{lem:bianchi}
$d_{CE}\widehat\Delta+\{S_{\Psi,\mathfrak{psl}},\widehat\Delta\}=0$ and $b_D^2=\{\widehat\Delta,-\}$; $\widehat\Delta=e^{\{\Psi,-\}}\widehat\Delta_{\rm nm}$ with $\widehat\Delta_{\rm nm}$ the breaking of $S_{{\rm nm},\mathfrak{psl}}$; $\mathrm{LIM}\,\widehat\Delta=0$.
\end{lemma}

\begin{proof}
The first two statements are the computation that gives $s^2=0$ from \eqref{eq:MC4} in [HK, thm:c8-obstruction], with the defect kept: $d_{CE}$ is a derivation of the cochain bracket, $d_{CE}^2=0$, and the Jacobi identity gives $\{S,\{S,F\}\}=\frac12\{\{S,S\},F\}$ for $S$ of degree $0$. The third is Lemma \ref{lem:canon}. At $D=4$ in the four-dimensional realisation $S_{{\rm nm},\mathfrak{psl}}$ solves \eqref{eq:MC4} [HK, thm:c2-master], the non-minimal pair being a contractible pair carrying the representation of Definition \ref{def:equiv}; $\mathrm{LIM}$ is that evaluation.
\end{proof}

\begin{theorem}[Classical breaking in Scheme D]
\label{thm:breakD}
Let $\mathscr L_0$ be the density of $S_{0,D}:=\int\mathrm{dvol}_D\big(-\frac14(F_A,F_A)+\frac12(\lambda,{\not d}_A\lambda)\big)$. In Scheme D:
\begin{itemize}
\item[(i)] $\widehat\Delta^{(0)}=0$ (gauge BRST exact, non-minimal sector included).
\item[(ii)] For every even generator with vector field $v$, $\widehat\Delta^{(1)}(x)=(4-D)\int\mathrm{dvol}_D\,\sigma_v\,\mathscr L_0$: zero along $P_{v_1}$, $M_\omega$, $R_r$; $2(4-D)S_{0,D}$ along $\Delta$; $-2(4-D)\int\mathrm{dvol}_D(v_1,x)\mathscr L_0$ along $K_{v_1}$. With $4-D=2\epsilon$ these are the BV forms of $\delta_DS_\epsilon=-\frac{2\epsilon}{g^2\mu^{2\epsilon}}\int d^Dx\,L$ and of the special conformal variation $\propto\frac{4\epsilon}{g^2\mu^{2\epsilon}}\int d^Dx\,x^\nu L$ of [DHKS07, (25), (27)] (up to the factor $2$ of $\Delta=2E$, the sign convention of $K$ and the factor $i$ of the weight $e^{iS}$ there).
\item[(iii)] Along $\xi\in\Sigma$: $\widehat\Delta^{(1)}(\xi)=\mathfrak F(\xi):=\frac12\int\mathrm{dvol}_D\big(\lambda,\rho([\xi.A,\lambda])\big)=-\frac12\int\mathrm{dvol}_D\big(\lambda,\rho([\Gamma(\xi,\lambda),\lambda])\big)$, cubic in $\lambda$; $\mathfrak F(\xi)$ vanishes in ten dimensions by [ESW, Assumption 3.1] and is Fierz-null.
\item[(iv)] Along $\eta\in\Sigma^*$, with $\mathfrak F_x(\zeta)$ the density of $\mathfrak F(\zeta)$:
\[
\widehat\Delta^{(1)}(\eta)=\sum_{\mu}\int\mathrm{dvol}_D\,x^\mu\,\mathfrak F_x(\rho(e_\mu)\eta)+(4-D)\,\mathcal N(\eta),\qquad \mathcal N(\eta):=-\tfrac12\int\mathrm{dvol}_D\big(\rho(F_A)\eta,\lambda\big),
\]
$\{e_\mu\}$ an orthonormal basis of $QDS$; the relative normalisation of the two terms is fixed, their common sign is that of $\xi.A=-\Gamma(\xi,\lambda)$.
\item[(v)] $\widehat\Delta^{(2)}(x,y)=0$ when $x$ or $y$ is even; for odd $x,y$, $\widehat\Delta^{(2)}(x,y)$ is linear in $\lambda^*$, Fierz-null, and vanishes in ten dimensions by [ESW, Lemma 3.6]; $\widehat\Delta^{(3)}$ on three odd arguments is linear in $c^*$ and vanishes in ten dimensions by [ESW, Lemma 3.7]; $\widehat\Delta^{(k\ge4)}=0$.
\end{itemize}
\end{theorem}

\begin{proof}
(i) The proof of $\{S,S\}=0$ in [HK, prop:c1-cme] uses gauge covariance of $F_A$ and ${\not d}_A$, the Leibniz rule and the Jacobi identity, all valid for formal $V^{\rm D}_{10}$ and any $D$; the non-minimal pair is contractible; $e^{\{\Psi,-\}}$ preserves the bracket.

(ii) By T3 the field coordinates transform as $\delta\phi=v(\phi)+w_\phi\sigma_v\phi$ plus the spin rotation by the antisymmetric part of $dv$, and antifields with $D-w_\phi$. For a covariant density $f$ of total weight $w_f$, $\int\mathrm{dvol}_D\,\delta f=\int\big(v(f)+w_f\sigma_vf\big)+R=(w_f-D)\int\sigma_vf+R$, where $R$ collects the terms with $d\sigma_v$. Antifield terms have $w_f=D$. For $S_{0,D}$, $w_f=4$: $F_A$ is a two-form and $L_vF_A$ has no $d\sigma_v$-term; the scalar kinetic term gives $R=-\int(\partial^2\sigma_v)(A|_E,A|_E)$ with $E=Q(4{-}D)S\oplus W^\perp$, zero because $\sigma_v$ is affine for every generator; the fermion kinetic term gives a multiple of $\int(\lambda,\rho(d\sigma_v)\lambda)$, zero because $d\sigma_v$ is constant and the bilinear form $(\zeta_1,\rho(u)\zeta_2)$ is symmetric while $\lambda$ is odd. The non-minimal and gauge-fixing parts are covered by Lemma \ref{lem:bianchi}: $\widehat\Delta=e^{\{\Psi,-\}}\widehat\Delta_{\rm nm}$ and $\widehat\Delta^{(1)}_{\rm nm}(x)$ contains no antifield.

(iii) In [HK, prop:c1-norm] and [ESW, Lemma 3.4] the identity $\{S,S^{(1)}(\xi)\}=0$ uses the Clifford relation and the Bianchi identity $d(dA)=0$ in the quadratic part, gauge covariance in the terms with antifields, and [ESW, Assumption 3.1] in the single term cubic in $\lambda$, the variation of $\frac12(\lambda,\rho([A,\lambda]))$ along $\xi.A$. Only the last is unavailable in $\mathrm{Cl}_D$.

(iv) Let $\epsilon(x)=\rho(x)\eta$. With $\xi.A=-\Gamma(\epsilon,\lambda)$, $\xi.\lambda=\frac12\rho(F_A)\epsilon$ for position-dependent $\epsilon$, the variation of the quadratic part is, after the integrations by parts of the constant case and the identities $\rho(v)\rho(F)=\rho(v\wedge F)+\rho(\iota_vF)$, $\sum_\mu\rho(e^\mu)\rho(\partial_\mu dA)=-\rho(d^*dA)$,
\[
\tfrac12\sum_\mu\int\mathrm{dvol}_D\big(\rho(e^\mu)\rho(F_A)\partial_\mu\epsilon,\lambda\big),\qquad \partial_\mu\epsilon=\rho(e_\mu)\eta .
\]
For $N,P$ in $QDS$, in $QDS\times E$ and in $E$, $\sum_\mu\rho(e^\mu)\rho(q(e^N\wedge e^P))\rho(e_\mu)=(D-4)$, $(D-2)$, $D$ times $\rho(q(e^N\wedge e^P))$. The term $-(\rho(A|_E)\eta,\lambda^*)$ of $S^{(1)}(\eta)$ contributes the $D$-independent coefficients $-2$ and $-4$ on the last two types, which cancel them at $D=4$ [HK, prop:c2-coeff]; the remainder is $\frac12(D-4)\int(\rho(F_A)\eta,\lambda)=(4-D)\mathcal N(\eta)$. The cubic term is the $x$-moment of (iii).

(v) Even--even components: $v\mapsto L_v+w\sigma_v$ is a representation of the conformal algebra of $QDS$ for every weight $w$, because $\sigma_{[v,v']}=v(\sigma_{v'})-v'(\sigma_v)$. Even--odd: the brackets of conformal Killing vectors with Killing spinors are $D$-independent (the computation of [HK, prop:c2-special] uses $d(v,-)=4(x,-)\wedge(v_1,-)$ and $\sigma_{K_{v_1}}=-2(v_1,x)$, valid for $x\in QDS$), and every term of $S^{(1)}(x)$ has total weight $D-\frac12\lambda_x$, $\lambda_x$ the $\mathrm{ad}_\Delta$-eigenvalue [HK, def:c12-gens], so the weight terms are $D$-independent. Odd--odd and arity three: [ESW, Lemmas 3.6, 3.7] use [ESW, Assumption 3.1] for the $\lambda^*$- and $c^*$-terms; the remaining terms use the Clifford relation and the Jacobi identity. Arity $\ge4$: [HK, prop:c2-tower], which uses only the field content of $S^{(2)}$.
\end{proof}

\begin{theorem}[Classical breaking in Scheme B]
\label{thm:breakB}
In Scheme B, with $\mathscr L_0$ the density of $S_{4S}+S_{\rm evan}$ without antifields:
\begin{itemize}
\item[(i)] $\widehat\Delta^{(0)}=0$: $S^{\rm B}_D$ is gauge invariant, $\lambda$ being in the adjoint and $d_A$ covariant on $QDS$.
\item[(ii)] $\widehat\Delta^{(1)}=0$ along $P_{v_1}$, $M_\omega$ and along $R_r$ for $r$ in the stabiliser $\mathfrak{so}(4)_J$ of $J$; along the nine remaining generators of $\mathfrak{su}(4)_R$, $\widehat\Delta^{(1)}(R_r)=\frac12\int\mathrm{dvol}_D\big(\lambda,[\rho_J(d|_{Q(D-4)S}),r.]\lambda\big)$, evanescent; along $\Delta$, $K_{v_1}$ the formula of Theorem \ref{thm:breakD} (ii) with the present $\mathscr L_0$.
\item[(iii)] Along $\xi$: $\widehat\Delta^{(1)}(\xi)=\mathfrak F^{\rm B}(\xi)+\mathcal E(\xi)$, with $\mathfrak F^{\rm B}$ the cubic term of Theorem \ref{thm:breakD} (iii) with $4S\oplus W^\perp_{\mathbb R}$-contractions and $\mathcal E(\xi):=\{S_{\rm evan},S^{(1)}_D(\xi)\}$, whose quadratic part
\[
\begin{aligned}
\mathcal E(\xi)|_2=\Big\{&-\tfrac14\int\mathrm{dvol}_D(dA,dA)\big|_{\rm hat}+\tfrac12\int\mathrm{dvol}_D\big(\lambda,\rho_J(d|_{Q(D-4)S})\lambda\big),\\
&-\int\mathrm{dvol}_D\big(\Gamma(\xi,\lambda)|_{4S\oplus W^\perp},A^*\big)+\tfrac12\int\mathrm{dvol}_D\big(\rho(dA|_{4S\oplus W^\perp})\xi,\lambda^*\big)\Big\}
\end{aligned}
\]
is nonzero and evanescent: the $Q(D{-}4)S$-components of the gauge field have no partner among the fermion components (field counting).
\item[(iv)] Along $\eta$: by T4 the ghost structures use $\rho(x|_{4S})\eta$, and $\widehat\Delta^{(1)}(\eta)=\sum_{\mu\in4S}\int\mathrm{dvol}_D\,x^\mu\mathfrak F^{\rm B}_x(\rho(e_\mu)\eta)+\mathcal E(\rho(x|_{4S})\eta)$: the four-dimensional sandwich $\sum_{\mu\in4S}\rho(e^\mu)\rho(q(e^N\wedge e^P))\rho(e_\mu)$ has the $D$-independent coefficients $0,2,4$, and the conformal breaking sits in $\mathcal E$.
\item[(v)] $\widehat\Delta^{(2)}(x,y)=0$ for $x,y$ even; along $(K_{v_1},\zeta)$ (whose bracket is realised with $x\in QDS$ on the left and $x|_{4S}$ on the right), along $(R_r,\zeta)$ with $r\notin\mathfrak{so}(4)_J$, and on odd pairs, $\widehat\Delta^{(2)}$ contains evanescent terms with $\rho_J(Q(D{-}4)S)$, $x|_{Q(D-4)S}$ and the $Q(D{-}4)S$-components of $\Gamma$, besides the Fierz-null terms of Theorem \ref{thm:breakD} (v).
\end{itemize}
\end{theorem}

\begin{proof}
(i) Gauge transformations act on $\lambda$ by $[c,-]$, commuting with $\rho$, $\rho_J$ and $J\otimes\tau$. (ii) Translations and rotations of $4S$ commute with $\gamma_5$, $\tau$ and $J$; $r\in\mathfrak{su}(4)$ commutes with $J\otimes\tau$ exactly on $\mathfrak{so}(4)_J$. The weight computation of Theorem \ref{thm:breakD} (ii) applies term by term, the evanescent kinetic term having weight $4$. (iii) The $4S\oplus W^\perp$-parts of the quadratic terms cancel as in four dimensions; the terms of $S_{\rm evan}$ vary under $\xi.A$, $\xi.\lambda$ (T4) into terms with one $Q(D{-}4)S$-derivative and no compensating term, the transformation of $A|_{Q(D-4)S}$ being zero by T4. (iv), (v) As (iii), with Theorem \ref{thm:breakD} (iv)--(v); with $\epsilon=\rho(x|_{4S})\eta$, $\partial_\mu\epsilon=\rho(e_\mu)\eta$ for $\mu\in4S$ and $0$ for $\mu\in Q(D{-}4)S$.
\end{proof}

\begin{corollary}[Which generators fail classically]
\label{cor:fail}
\mbox{}\begin{itemize}
\item Exact in both schemes: gauge BRST, $P_{v_1}$, $M_\omega$; exact in Scheme D, and on $\mathfrak{so}(4)_J$ in Scheme B: $R_r$.
\item Broken by the dimensionful coupling, $\widehat\Delta=(4-D)\times$(weighted action density): $\Delta$, $K_{v_1}$, and the part $(4-D)\mathcal N(\eta)$ along $\eta$. The coupling has mass dimension $D-4$ in the pairing: $\int\mathrm{dvol}_D(F_A,F_A)$ has weight $4-D$.
\item Broken through Fierz identities: $\xi$, $\eta$ and the odd--odd closure (arity two and three), in both schemes.
\item Broken through field counting: in Scheme B along $\xi$, $\eta$ ($D+6$ bosonic against sixteen fermionic components); in Scheme D the counting $D+(4-D)+6=10$ is exact and the free part of $\{S_D,S^{(1)}_D(\xi)\}$ vanishes identically.
\end{itemize}
\end{corollary}

\begin{remark}[Two-component form of the Fierz breaking]
For $\xi\in S_+\otimes W$, with $\Gamma:=\Gamma(\xi,\lambda)$, $\Gamma^m_V=i\bar\lambda_{\dot\alpha i}\bar\sigma^{m\dot\alpha\alpha}\xi_\alpha{}^i$, $\Gamma_{kl}=2\varepsilon_{jikl}\xi_\alpha{}^j\varepsilon^{\alpha\beta}\lambda_\beta{}^i$, $\Gamma^{kl}=2(\xi^k\varepsilon\lambda^l-\xi^l\varepsilon\lambda^k)$ [HK, def:c8-twocomp, rem:c12-twocomp],
\[
\begin{aligned}
\mathfrak F(\xi)=-\tfrac12\int d^Dx\Big\{&i\big(\lambda_\alpha{}^i,[\Gamma^m_V,\bar\lambda_{\dot\alpha i}]\big)\bar\sigma^{m\dot\alpha\alpha}+\varepsilon^{\alpha\beta}\big(\lambda_\alpha{}^i,[\Gamma_{ij},\lambda_\beta{}^j]\big)\\
&+i\big(\bar\lambda_{\dot\alpha i},[\Gamma^m_V,\lambda_\alpha{}^i]\big)\bar\sigma^{m\dot\alpha\alpha}-\varepsilon^{\dot\alpha\dot\beta}\big(\bar\lambda_{\dot\alpha i},[\Gamma^{ij},\bar\lambda_{\dot\beta j}]\big)\Big\},
\end{aligned}
\]
by $(\lambda,\rho(v)\lambda')=\lambda_\alpha{}^i(\rho(v)\lambda')^\alpha{}_i-\bar\lambda_{\dot\alpha i}(\rho(v)\lambda')^{\dot\alpha i}$ and [HK, eq:c8-pairing]. In four dimensions it vanishes by $\sigma^m_{\alpha\dot\alpha}\sigma^m_{\beta\dot\beta}=-2\varepsilon_{\alpha\beta}\varepsilon_{\dot\alpha\dot\beta}$ and the Schouten identity; it is the $\mathcal N=4$ counterpart of $\Delta_{\rm gauge}=\frac i2gf_{abc}(\bar\epsilon\gamma^\mu\tilde g_a)(\bar{\tilde g}_b\gamma_\mu\tilde g_c)$ of [S05, (4.4b)]. In Weinberg's form replace $\rho(v)$ by $\gamma_W^\mu v_\mu$ on the Dirac column.
\end{remark}

\begin{remark}[Two-component forms of the other breakings]
\label{rem:ws2}
With [HK, eq:c8-pairing] and $(\rho(w)\eta)_\alpha{}^k=-w^{kj}\varepsilon_{\alpha\beta}\eta^\beta{}_j$, $(\rho(w)\eta)_{\dot\beta k}=-w_{kj}\varepsilon_{\dot\beta\dot\alpha}\bar\eta^{\dot\alpha j}$:
\begin{itemize}
\item $\mathcal N(\eta)=-\frac12\int d^Dx\,\big[\lambda_\alpha{}^i(\rho(F_A)\eta)^\alpha{}_i-\bar\lambda_{\dot\alpha i}(\rho(F_A)\eta)^{\dot\alpha i}\big]$ (colour pairing understood), where the part of $\rho(F_A)\eta$ along $QDS\wedge QDS$ is
\[
(\rho(F_A)\eta)^\alpha{}_i\supset\tfrac12F_{mn}\,\eta^\beta{}_i(\sigma^m\bar\sigma^n)_\beta{}^\alpha,\qquad (\rho(F_A)\eta)^{\dot\alpha i}\supset-\tfrac12F_{mn}(\bar\sigma^m\sigma^n)^{\dot\alpha}{}_{\dot\beta}\,\bar\eta^{\dot\beta i},
\]
the part along $QDS\wedge E$ is $\sum_m\rho(e^m)\rho(D_mA|_E)\eta$, i.e. $(\cdot)^\alpha{}_i\supset-i\,D_mA_{ij}\,\varepsilon_{\dot\alpha\dot\beta}\bar\eta^{\dot\beta j}\bar\sigma^{m\dot\alpha\alpha}$ and $(\cdot)^{\dot\alpha i}\supset i\,\bar\sigma^{m\dot\alpha\alpha}D_mA^{ij}\varepsilon_{\alpha\beta}\eta^\beta{}_j$, and the part along $E\wedge E$ is $\frac12\sum_{w\neq w'}[A_w,A_{w'}]\rho(e^w)\rho(e^{w'})\eta$ composed from the same formulas; in Scheme D the $\epsilon$-scalar components $A_{\tilde m}$ enter through $\rho(e^{\tilde m})$ with $\tilde m\in Q(4{-}D)S$ formal.
\item Scheme B, $R$-symmetry: for $r\in\mathfrak{su}(4)$ acting by $\lambda_\alpha{}^i\mapsto r^i{}_k\lambda_\alpha{}^k$, with $r_{ik}:=\delta_{ij}r^j{}_k$,
\[
\widehat\Delta^{(1)}(R_r)=\tfrac12\int d^Dx\Big\{\varepsilon^{\alpha\beta}(r_{ik}+r_{ki})\big(\lambda_\alpha{}^i,\tau^{\hat m}D_{\hat m}\lambda_\beta{}^k\big)-\varepsilon^{\dot\alpha\dot\beta}(r^{ik}+r^{ki})\big(\bar\lambda_{\dot\alpha i},\tau^{\hat m}D_{\hat m}\bar\lambda_{\dot\beta k}\big)\Big\},
\]
zero exactly on $\mathfrak{so}(4)_J$ ($r_{ik}$ antisymmetric).
\item Scheme B, field counting: $\mathcal E(\xi)|_2$ is the bracket of the evanescent quadratic terms $-\frac18(F_{\hat mN},F_{\hat mN})-\frac1{16}(\partial_{\hat m}A^{ij},\partial_{\hat m}A_{ij})+\frac12\{\delta_{ij}\varepsilon^{\alpha\beta}(\lambda_\alpha{}^i,\tau^{\hat m}\partial_{\hat m}\lambda_\beta{}^j)-\delta^{ij}\varepsilon^{\dot\alpha\dot\beta}(\bar\lambda_{\dot\alpha i},\tau^{\hat m}\partial_{\hat m}\bar\lambda_{\dot\beta j})\}$ ($\hat m\in Q(D{-}4)S$, $N\in4S\oplus Q(D{-}4)S$) with the linear terms of $S^{(1)}(\xi)$ in the two-component form of [HK, rem:c12-twocomp]: $T_1=-i\int(\bar\lambda_{\dot\alpha i}\bar\sigma^{m\dot\alpha\alpha}\xi_\alpha{}^i,A^*_m)$, $T_2=-\frac12\int\varepsilon_{jikl}(\xi_\alpha{}^j\varepsilon^{\alpha\beta}\lambda_\beta{}^i,A^{*kl})$, $T_3=-\frac14\int((\partial_mA_n-\partial_nA_m)\sigma^m_{\alpha\dot\alpha}\bar\sigma^{n\dot\alpha\beta}\xi_\beta{}^i,\lambda^{*\alpha}{}_i)$, $T_4=-\frac i2\int(\varepsilon^{\alpha\beta}\xi_\beta{}^j\sigma^m_{\alpha\dot\alpha}\partial_mA_{ij},\bar\lambda^{*\dot\alpha i})$, $m,n\in4S$.
\item In the non-minimal Feynman gauge the ghost line is $\langle c\,\bar c\rangle\propto(-,-)^{-1}|p|^{-2}$, replacing the line $\langle c\,A^*_m\rangle$ of [HK, prop:c8-twocomp]; the lines $\langle A_mA_n\rangle$, $\langle A_{ij}A_{kl}\rangle$, $\langle\lambda_\alpha{}^i\bar\lambda_{\dot\beta j}\rangle$ are those of [HK, prop:c8-twocomp] with $|p|^{-2}$ and $D$-dimensional $p$.
\end{itemize}
\end{remark}

